Esta seção apresenta uma revisão de trabalhos relacionados que abordam SLAM visual, robótica móvel de baixo custo, sensores RGB-D e integração de sistemas robóticos em contextos próximos ao proposto. A análise concentra-se no desafio de executar percepção visual e mapeamento em hardware limitado, mantendo a coerência entre sensores, odometria e software de integração.

\section{SLAM Visual e Hardware Embarcado}

A implementação de SLAM em hardware embarcado de baixo custo impõe desafios severos. \citeonline{gatesichapakorn2019rosbasedautonomous} propuseram uma implementação de navegação autônoma utilizando o \textit{stack} de navegação do ROS, realizando experimentos comparativos entre um Raspberry Pi 3, empregado com LiDAR 2D, e um Intel NUC, empregado com LiDAR e câmera RGB-D. Os resultados demonstraram que, embora o Raspberry Pi fosse capaz de sustentar cenários mais simples, o processamento de dados visuais exigiu hardware mais robusto para uma fusão sensorial eficiente. Esse resultado evidencia o gargalo de processamento enfrentado por plataformas de baixo custo.

No contexto de processamento inteiramente embarcado (\textit{onboard}), \citeonline{canovas2021onboarddynamicrgbd} propuseram uma abordagem de SLAM visual RGB-D resiliente a ambientes dinâmicos. Embora tenham alcançado resultados robustos na filtragem de objetos móveis, o estudo destaca o alto custo computacional de manter a estimativa de pose e o mapeamento denso simultaneamente na CPU do robô. Essa limitação reforça a necessidade de avaliar cuidadosamente frequência, resolução, parâmetros do algoritmo e estabilidade térmica quando o processamento ocorre no próprio robô.

Para viabilizar o vSLAM em hardware restrito, a otimização de parâmetros é crucial. \citeonline{nagarajan2024multiobjectiveoptimizationrtabmap} investigaram especificamente o algoritmo RTAB-Map, também selecionado neste trabalho, e propuseram o uso de Algoritmos Genéticos Multiobjetivo (MOGA) para ajustar automaticamente seus parâmetros. O estudo demonstrou que a configuração padrão do software nem sempre oferece o melhor desempenho em cenários \textit{indoor}, evidenciando que o ajuste fino (\textit{tuning}) é uma etapa indispensável para equilibrar qualidade do mapa e consumo de recursos.

\section{Sensores RGB-D, Odometria e Fusão Sensorial}

O uso de câmeras RGB-D em robótica móvel de baixo custo é atrativo porque fornece, simultaneamente, cor e profundidade, permitindo a construção de mapas 3D sem sensores LiDAR de maior custo. Entretanto, essa escolha introduz limitações práticas, como alcance reduzido, sensibilidade à iluminação, campo de visão restrito e elevado volume de dados. \citeonline{kamarudin2015improvingperformance2d} demonstram que sensores RGB-D podem apresentar desempenho inferior a LiDARs em determinados cenários, especialmente quando o ambiente possui poucas características visuais ou exige alcance maior.

Além da percepção visual, a qualidade da odometria é determinante para a estabilidade do SLAM. Sistemas RGB-D podem perder rastreamento quando há movimentos bruscos, superfícies pouco texturizadas ou inconsistência entre a movimentação real e a estimativa inicial de pose. Nesse contexto, \citeonline{qayyum2017imuaidedrgbd} discutem o uso de informação inercial para auxiliar sistemas RGB-D, indicando que a fusão entre IMU e visão pode contribuir para maior robustez em situações de movimento rápido ou ruído visual.

Esses trabalhos fundamentam a estratégia adotada nesta pesquisa: tratar a câmera RGB-D como sensor principal de mapeamento, mas avaliar também a coerência da odometria das rodas e da IMU como fatores que influenciam diretamente a estabilidade do RTAB-Map no robô físico.

\section{Integração entre Percepção Visual e Navegação}

Um desafio recorrente no uso de vSLAM é a tradução dos dados 3D para estruturas úteis à navegação. \citeonline{zhang2023rgbdnavigation2d} apresentaram o \textit{framework} ``RGBD Navigation'', que propõe converter nuvens de pontos 3D e mapas de profundidade em varreduras a laser 2D virtuais. Isso permite a integração direta com pilhas de navegação baseadas em mapas de ocupação, oferecendo um caminho para aproximar percepção visual rica e planejadores de trajetória 2D.

Validando o uso dessas ferramentas em ambiente controlado, \citeonline{kulkarni2021visualslamcombined} utilizaram a biblioteca RTAB-Map integrada ao ROS e ao simulador Gazebo para navegação autônoma combinada com detecção de objetos. Embora a simulação não faça parte do eixo de validação desta versão do trabalho, esse estudo é relevante por demonstrar a integração entre RTAB-Map, ROS e tarefas de navegação em uma arquitetura robótica completa.

Funcionalidades de navegação autônoma completa, como planejamento de trajetórias e desvio de obstáculos com Nav2, permanecem como continuidade possível para esta pesquisa. A etapa atual, porém, concentra-se primeiro em consolidar o mapeamento, a localização e a estabilidade do pipeline sensorial no robô físico.

\section{Contexto Institucional e Síntese}

No cenário nacional e institucional, \citeonline{lemos2019plataformaroboticamovel} desenvolveram uma plataforma móvel baseada em Kinect e RTAB-Map, evidenciando as limitações de processadores de baixo consumo para processar nuvens de pontos em tempo real. Mais recentemente, \citeonline{araujo2023atmosbotpropostamodelagem}, em trabalho desenvolvido na mesma instituição deste projeto, propôs o \textit{AtmosBot}, uma modelagem simulada focada na navegação autônoma utilizando sensores LiDAR 2D.

A presente pesquisa diferencia-se desses trabalhos ao priorizar a integração física de uma plataforma robótica de baixo custo com percepção RGB-D, controle embarcado, IMU, odometria e ROS 2. Em relação ao trabalho institucional baseado em LiDAR 2D, o avanço está na transição para percepção visual 3D em hardware real. Em relação aos estudos de SLAM visual em hardware mais robusto, a contribuição está em documentar os limites e ajustes necessários para operar com Raspberry Pi e componentes acessíveis.

A Tabela \ref{tab:trabalhos_correlatos} sintetiza as características dos trabalhos discutidos em relação à proposta atual.

\begin{table}[htb]
\centering
\caption{Comparativo entre trabalhos correlatos e a proposta atual.}
\label{tab:trabalhos_correlatos}
\resizebox{\textwidth}{!}{%
\begin{tabular}{|p{3.5cm}|p{2.5cm}|p{3cm}|p{5cm}|}
\hline
\textbf{Autor/Ano} & \textbf{Sensor Principal} & \textbf{Contexto} & \textbf{Contribuição Principal} \\ \hline
\citeonline{araujo2023atmosbotpropostamodelagem} & LiDAR 2D & Simulação & Navegação e modelagem robótica em Gazebo \\ \hline
\citeonline{gatesichapakorn2019rosbasedautonomous} & LiDAR + RGB-D & Robô físico & Comparação de desempenho entre sensores e plataformas \\ \hline
\citeonline{lemos2019plataformaroboticamovel} & RGB-D & Robô físico & Limitações de hardware em plataforma móvel com RTAB-Map \\ \hline
\citeonline{canovas2021onboarddynamicrgbd} & RGB-D & Robô físico & SLAM visual embarcado em ambientes dinâmicos \\ \hline
\citeonline{nagarajan2024multiobjectiveoptimizationrtabmap} & RGB-D & Avaliação de algoritmo & Otimização de parâmetros do RTAB-Map \\ \hline
\citeonline{zhang2023rgbdnavigation2d} & RGB-D & Integração de software & Conversão de percepção RGB-D para navegação 2D \\ \hline
\citeonline{kulkarni2021visualslamcombined} & RGB-D & Simulação ROS/Gazebo & Integração entre RTAB-Map, ROS e navegação \\ \hline
\textbf{Proposta Atual} & \textbf{RGB-D + IMU} & \textbf{Robô físico de baixo custo} & \textbf{Integração e avaliação de SLAM visual embarcado em Raspberry Pi} \\ \hline
\end{tabular}
}
\fonte{Elaborado pelo autor.}
\end{table}
